\section{Admission precedes interpretation in a causal client}
\label{exec17:i2-admission}

The three surfaces become useful to a consumer only when the consumer can
decide what kind of reply it has received. Consider a scheduled lending
report that asks for a conditional adjustment expression, saves the reply,
and later displays it through an agent. The numerical estimate is correctly
absent. If the final client converts that absence to zero, the report now
asserts a null monetary effect that no component calculated. If it displays
only a successful process status, the unsupported and refused cases look
equally favorable. The missing product component is therefore a reviewed
client interpretation contract, with an explicit boundary between receiving
bytes, admitting a response, and deciding how that response may be used.

A proposed client session should retain its launch configuration and admitted
server identity, the negotiated protocol where applicable, the operation
catalogue and schemas, and the exact scientific request. The catalogue is
evidence about available operations in that session. It is not a warranty
that every graph class or estimator mentioned in a description is supported.
Admission needs the operation, method identifier and declared scientific
domain together. In the inspected local implementation, the outer MCP tool
name must agree with the operation inside the canonical request; discovery
is absent from the available mapping. These are actual source boundaries
\citep{exec13_i2_examples}. The more detailed consumer decision procedure
below is a proposed extension, evaluated by a separate finite reference
rather than installed as a new production client.

\subsection{A narrow decoding profile with explicit failure states}

First establish the transport boundary. A CLI consumer has process exit,
stdout and stderr as separate observations. A protocol consumer has a
correlated protocol reply or a transport failure before any causal response
exists. Diagnostic output should be retained separately from the object to
be decoded. Stripping an arbitrary prefix until some braces can be parsed
would hide a corrupted channel and could select the wrong object from a
mixed output stream. The proposed local profile accepts one complete JSON
value at the expected boundary, subject to its byte allowance. It rejects a
banner or trailing unrelated text. This is a deliberate consumer restriction,
not a claim that every possible transport framing uses one JSON document.

The profile also rejects duplicate object keys, nonfinite numeric tokens and
numeric overflow that becomes nonfinite after parsing. Duplicate keys matter
because two parsers can choose different values for the same field. A reply
containing both a completed and a failed status under the same key has no
single admitted meaning in this profile. Permitting a nonfinite number as a
point estimate would also bypass the explicit representation of unbounded
sets. The decoder consequently checks both the lexical input and the parsed
numeric tree. A finite byte limit reduces the amount of accepted input but
does not by itself bound all parsing memory or nesting depth; a deployment
must add and measure those resource boundaries.

After decoding, validate the full request and response schemas and the
semantic relation between them. The local semantic validator checks request
and operation identity, the context digest, units, result state, interval
topology and uncertainty interpretation. For example, a result marked
uncomputed cannot carry a numerical estimate, and a sensitivity result has
no sampling-confidence level. These are testable declared relationships.
They do not establish that the graph or supplied covariance is true
\citep{exec13_i2_examples}. The proposed client should keep schema admission
and semantic admission as separate diagnostic steps so that a missing field
does not become a scientific claim of nonidentification.

The finite reference applies the existing semantic validator to eleven
preserved library/CLI fixtures and classifies their meanings. It separately
rejects duplicate keys, a nonfinite token, overflow, trailing text, a banner,
an excessive byte count, a wrong request identifier, a wrong operation,
a wrong context, inconsistent set topology and an unsupported response
carrying a numerical result. A twelfth adverse case compares conflicting
structured and text representations. The exercise uses local stored
responses and deliberate copies; it makes no network call and does not
constitute a new MCP transport trial. Its scope is the proposed semantic
consumer, not complete schema, SDK, parser-resource or deployed-client
qualification \citep{exec17_integration_review_reference}.

\subsection{Two representations must describe one admitted response}

The inspected MCP adapter returns structured content and a serialized text
copy of the same canonical response. A proposed client that consumes both
should decode them separately and compare their canonical finite JSON
representations before interpreting either. Object ordering and whitespace
can vary without changing the declared fields, while ordered arrays remain
ordered. The reference uses the project's explicit serialization convention;
it does not claim a universal JSON-equivalence theorem. In particular, the
chosen identity need not treat every numerically equivalent spelling as
interchangeable. Any normalization beyond that convention must be justified
as part of the admitted profile.

A disagreement in qualification is consequential even when the estimated
coefficient agrees. One copy might say execution-only while another says
conditional inference. A disagreement in interval topology or context can
be still more serious. The proposed response to any such conflict is an
integrity failure with both originals retained. Choosing whichever copy gives
a convenient answer would make the consumer an unreviewed arbiter of
scientific meaning. If a permitted protocol profile supplies only one
representation, the client can use that representation after its own checks;
absence of a second copy is not automatically a conflict. The local reference
checks agreement under changed whitespace and rejection after changing a
qualification field \citep{exec17_integration_review_reference}.

This decision also specifies when a language model should see the result.
A deterministic admission layer should first establish which object exists
and what fields it contains. The model can then explain a supported typed
object with links to its assumptions and evidence. Asking the model to repair
a malformed response before preserving the original creates an additional
unmeasured transformation. A corrected result must instead have its own
identity and a recorded reason for the change. This is a design obligation
for the proposed client, not an assertion that the current local service
already implements a complete user interface.

\subsection{The decision table for a lending report}

Once the response is admitted, its interpretation follows the result type.
The following branches are deliberately separate because their remedies
require different information.
\begin{center}
\begin{tabularx}{\textwidth}{P{.24\textwidth}X}
\toprule
Admitted outcome & Proposed consumer behavior \\
\midrule
Execution failed or rejected & Retain error and request identity; publish no scientific answer. \\
Scientific unsupported & Show the unsupported domain or method; preserve the requested target. \\
Scientific refused & Show the contradictory or unmet premise; do not retry it as a numerical failure. \\
Identification expression & Display the expression and assumptions; do not manufacture an effect estimate. \\
Effect set & Preserve every component and its uncertainty interpretation, including empty or unbounded results. \\
Job reply & Read the payload's job state; expose a report only through its admitted completed result reference. \\
\bottomrule
\end{tabularx}
\end{center}
This table is a display and routing rule. It cannot promote execution-only
evidence to scientific qualification. The preserved fixtures demonstrate
conditional expressions, unsupported sampling, contradicted premises,
finite/empty/whole-line/disconnected sets and context rejection. The proposed
consumer keeps these outcomes distinct
\citep{exec13_i2_examples,exec17_integration_review_reference}.

For an effect set, the display should show the monetary coordinate, sign,
horizon and the construction of the set. An empty sensitivity set means that
the supplied restrictions have no joint solution in that coordinate; it does
not mean a zero effect or zero uncertainty. A disconnected set cannot be
displayed as its central gap, and an unknown set cannot be converted to the
whole line merely because both are inconvenient to plot. The mathematical
distinctions are developed in the interchange chapter. The additional client
obligation here is to make each distinction survive storage, filtering,
sorting, plotting and export. For example, dropping rows with a null scalar
estimate would remove correctly computed identification expressions and
unbounded effect sets from the report denominator.

A completed job-status query is another important branch. The query can
finish successfully while the job payload says running. The local reference
constructs two semantic job replies with the same completed query status
and payloads running and completed, and obtains different display states.
These are authored semantic fixtures, not submitted report jobs. A full
client must additionally validate the job request schema, result-reference
presence, role, size, custody and bytes before opening a completed report.
It must then inspect the report's individual declarations. The inspected
report builder summarizes qualification; it does not supply a new causal
review of its inputs \citep{exec13_i2_examples}.

\subsection{Lost replies, stale observations and consumer state}

The client needs its own durable record of an attempt. A timeout after
submission does not establish that the server did no work. If a stable
idempotency identity was supplied, the consumer can use the permitted replay
or status path and retain the original request. Changing the scientific
context or budget while keeping the same identity requests a different
computation and should follow the conflict branch. The current implementation
binds the selected execution environment as well as the request, so replay
after a source change can be rejected even when the business sentence is
unchanged. That rejection is evidence about compatibility, not a failure to
understand the user's wording \citep{exec13_i2_examples}.

Polling introduces observation order. Suppose one status request begins
before completion but its response arrives after a later request that saw
the completed state. A naive display that overwrites its state with every
arriving reply can regress from completed to running. The proposed client
should associate observations with the exact job and request identity and
apply a declared terminal-state rule. Under an immutable job identity, a
terminal observation is not replaced by a later-arriving nonterminal
observation. Conflicting terminal states require reconciliation with the
authoritative job record; choosing the more favorable state is not a repair.
This is a proposed consumer rule and an adverse case for a future real
transport test, not a claim that the finite semantic reference simulates
network ordering or distributed consistency.

Cancellation has the same distinction between a request and an observation.
Recording that the user requested cancellation is compatible with a final
completed job when publication won the race. The client should preserve both
facts and label the resulting report accordingly. It must not equate closing
a window, cancelling a protocol call, requesting job cancellation and undoing
an external financial action. The actual service exposes report jobs and
does not perform those external actions. A consumer that later adds them
needs a separate authoritative reconciliation contract, not just another
status string.

\subsection{Capability negotiation determines a feasible workflow}

A proposed workflow planner should admit a chain of operations only when
each step's typed output can satisfy the next step's input domain. A
back-door expression is not a moment-summary artifact; a graph suggestion
is not a validated causal DAG; a successful timestamp diagnostic is not an
assignment mechanism. Therefore operation-name agreement alone cannot
establish composability. The current common runtime is useful precisely
because it exposes those narrow boundaries. The consumer should show the
unsupported bridge explicitly and retain an analyst path when no admitted
bridge exists.

For each candidate route, the planner should compare the declared target,
graph class, observation law, input artifact role, supported method and
result interpretation. If any component is absent, it can propose an
alternative route or ask for the missing scientific choice. It must not
silently replace a dependent panel by independent rows to gain admission.
For a source or protocol upgrade, the consumer's recorded capability snapshot
helps identify which prior route may have changed. Re-admission requires a
semantic comparison and selected adverse fixtures; a changed version string
neither proves compatibility nor automatically invalidates every earlier
scientific result.

The practical prototype deliverable is thus a small deterministic consumer
with a visible admission trace, followed by an optional explanation layer.
Its evaluation should pair each supported result with an appropriate
negative case and inspect the stored and exported objects, not merely the
screen caption. The finite reference already makes specific decoding and
interpretation errors reproducible. Completing an actual client would add
full schema validation, bounded parser resources, live transport ordering,
artifact retrieval, custody policy and launch-environment trials. These are
known engineering steps justified by the proposal. They remain separate
from expensive bank experiments, which should follow only after the
prototype demonstrates useful preservation of the scientific question.
\label{exec17:i2-admission-end}
